\bookchapter{第二章：整理汝之數字之室——文件管理}

\chapterintro{手機用久則亂，善習始於文件之管理。猶汝之室須定時整理，汝之數字空間亦須打理。井然有序之手機，尋物既便，運行亦暢。}

\booksection{一、三層存儲之體：汝之數字倉}

\begin{enumerate}
  \item \textbf{機身存儲}：核心之區，速度迅捷，系統與App皆在於此。此乃汝手機至寶之空間，宜精打細算。
  \item \textbf{SD卡（若支持）}：擴展之衣櫃，速度稍緩，宜於存放電影、照片等大文件。注：今諸多旗艦機已不支持SD卡。
  \item \textbf{雲存儲}：保險之櫃兼同步之網盤。防丟失，亦能於多設備間同步。iCloud、百度網盤、OneDrive皆常用之選。
\end{enumerate}

\booksection{二、識“文件管理器”}

安卓一般有 [文件] 或 [文件管理] App，iOS乃自帶之 [文件] App。乃汝檢視手機諸文件之“資源管理器”也。

\booksubsection{常見文件夾之義：}

\begin{itemize}
  \item \textbf{DCIM}：相機所攝之照片與視頻。
  \item \textbf{Pictures}：他圖，如截圖、所存之圖。
  \item \textbf{Download}：瀏覽器所下之文件。
  \item \textbf{Music}：音樂文件。
  \item \textbf{Movies}：視頻文件。
  \item \textbf{Documents}：文檔文件。
  \item \textbf{Android}：安卓系統之文件夾，內爲各App之數據，勿手動刪之。
\end{itemize}

\booksection{三、文件系統之結構：深入知之}

安卓本於Linux，文件系統有層級之構。雖常人不必深究，然知其基本之義亦有所助。

\booksubsection{存儲分區：}

\begin{itemize}
  \item \textbf{/system}：系統分區，存操作系統之文件。常人無權訪問，須root之權。
  \item \textbf{/data}：數據分區，存用戶數據與應用數據。此乃汝日常所用之空間。
  \item \textbf{/sdcard 或 /storage/emulated/0}：用戶可訪問之存儲空間，汝之照片、下載文件皆在於此。
\end{itemize}

\booksubsection{App數據之存儲：}

\begin{itemize}
  \item \textbf{內部存儲}：App之私有空間，他App不能訪問。卸載App時則刪之。
  \item \textbf{外部存儲}：App於用戶可訪問處所存之文件，如照片、下載之物。卸載App時未必刪之。
\end{itemize}

\booksection{四、數據之備份與恢復：爲汝數據上保險}

數據無價，備份乃必養之習。

\booksubsection{備份之策：}

\begin{enumerate}
  \item \textbf{雲備份}：自動而便，然須網絡與雲空間。iCloud、華爲雲、小米雲等皆支持。
  \item \textbf{本地備份}：以數據線連電腦，用官方之具備份。速快、空間大，然須手動操作。
  \item \textbf{雙重備份}：雲備份加本地備份，至爲安全。
\end{enumerate}

\booksubsection{備份何物？}

\begin{itemize}
  \item \textbf{照片與視頻}：最佔空間，亦最珍貴。
  \item \textbf{通訊錄}：存於雲賬號中最爲便。
  \item \textbf{微信聊天記錄}：用微信自帶之備份功能。
  \item \textbf{重要文檔}：存于云盤或電腦上。
\end{itemize}

\booksection{五、文件之安：護汝之隱私}

\begin{itemize}
  \item \textbf{敏感文件加密}：諸手機有"安全文件夾"或"私密空間"之功能，可將敏感文件納其中，須密碼或指紋方能訪問。
  \item \textbf{刪除非銷燬}：常之刪除僅標記空間可覆，數據猶在。欲徹底刪除，須用專門之安全擦除之具，或多次覆寫。
  \item \textbf{云端之安}：擇信譽佳之雲服務商，開雙重之認證，定期查登錄之設備。
\end{itemize}

\booksection{六、空間不足乎？來一番大掃除}

\begin{itemize}
  \item \textbf{清理緩存}：App之臨時記憶，清之則騰地。於設置 > 存儲 中可清理。
  \item \textbf{尋大文件}：用文件管理器按大小排序，揪出佔空間之"元兇"。
  \item \textbf{刪殘留}：卸載App後，時有垃圾之文件殘留，宜手動清之。
  \item \textbf{清理相冊}：刪重複之照、模糊之照、不需之截圖。
  \item \textbf{清理微信}：微信乃空間之大戶，定期清理聊天記錄與緩存甚爲要。
\end{itemize}

\booksection{七、常見文件格式全解析}

手機中諸文件後綴名令人頭眩，明乎其義可助汝善管文件：

\begin{center}
\begin{tabularx}{\linewidth}{|X|X|X|X|}
\hline
\textbf{格式} & \textbf{類} & \textbf{類} & \textbf{特性} \tabularnewline
\hline
JPG/JPEG & 圖 & 圖 & 相冊/任意圖覽器 \tabularnewline
相冊/任意圖覽器 & 圖 & 圖 & 相冊/任意圖覽器 \tabularnewline
相冊/任意圖覽器 & 音 & 音 & 有損壓縮，體小，兼容性極佳。最通用之音頻格式。 \tabularnewline
任意音樂播放器 & 音 & 音 & 無損壓縮，音質與CD同，體比MP3大三至五倍。發燒友之首擇。 \tabularnewline
音樂播放器/專業播放App & 音 & 音 & 無壓縮之原始音頻，音質至佳而體巨。專業錄音常用。 \tabularnewline
音樂播放器 & 音 & 音 & AAC之編碼，同碼率下音質勝於MP3。蘋果生態常用。 \tabularnewline
Apple Music/音樂播放器 & MKV & 視頻 & 最通用之視頻封裝格式，兼容性極佳。手機錄視頻之默認格式。 \tabularnewline
相冊/任意視頻播放器 & MKV & 視頻 & 開源之封裝格式，支持多音軌、多字幕。電影資源常用。 \tabularnewline
VLC/nPlayer & PDF & 文檔 & Excel電子表格。XLSX乃壓縮格式。 \tabularnewline
WPS/Office & PPT/PPTX & 演示 & Excel電子表格。XLSX乃壓縮格式。 \tabularnewline
WPS/Office & XLS/XLSX & 表格 & Excel電子表格。XLSX乃壓縮格式。 \tabularnewline
WPS/Office & PDF & 文檔 & 跨平臺之格式，排版固定不可編輯。合同、論文、電子書常用。 \tabularnewline
Adobe Reader/WPS & RAR & 壓縮 & 最通用之壓縮格式，跨平臺兼容性佳。 \tabularnewline
系統自帶/解壓App & RAR & 壓縮 & 壓縮率高於ZIP，然須專用軟件解壓。 \tabularnewline
解壓App & ISO & 鏡像 & 光盤鏡像之文件，完整複製光盤之內容。系統安裝包常用。 \tabularnewline
解壓App/虛擬光驅 & HTML & 網頁 & 網頁之文件，用瀏覽器開之。所存之網頁常用此格式。 \tabularnewline
\hline
\end{tabularx}
\end{center}

\begin{tcbpractice}{練一練：整理"下載"文件夾}
練一練：整理“下載”文件夾

\end{tcbpractice}

\booksection{八、圖與文件系統之管理：使資產井然有序}

上節吾等習於“識文件夾”，然識之猶未足也。本節所講者，\textbf{何以主動建一套可久行之管理規範}，使汝之文件隨時可尋、不失、不爲人窺。

\booksubsection{文件之分類與命名之規}

亂之源乃無規。建規則實爲易事，要在於\textbf{堅持}。

\begin{itemize}
  \item \textbf{"時間+事件"命名之法}：文件夾之名帶時間戳與事件之述，便於日後回憶與檢索。如\texttt{202606\_旅行票據}、\texttt{2025年終總結}、\texttt{2024孩子生日}。如此既見時間之線，亦可一目知其內容。
  \item \textbf{三級分類之體}：
  \begin{itemize}
    \item \textbf{一級}：按文件之類（文檔 / 圖 / 視頻 / 音）
    \item \textbf{二級}：按使用之場（工作 / 生活 / 學習）
    \item \textbf{三級}：按具體之事件或項目（如"2026Q2\_項目A"）
  \end{itemize}
  \item \textbf{月度清理之計劃}：每月固定一日（建議發薪之日或月初一日），對"下載"目錄集中整理：\textbf{有用者歸檔、無用者刪、珍貴者轉存}。一年十二次小整理，勝於積一年而大掃除。
\end{itemize}

\begin{tcbpractice}{本月任務清單}
本月任務清單

\begin{enumerate}
  \item 于文件管理器中建\texttt{2026Q2\_項目A}形之本月工作文件夾；
  \item 形之本月工作文件夾；
  \item 爲今年之家庭照建\texttt{2025家庭}、\texttt{2026家庭}此形之年度文件夾。
\end{enumerate}

\end{tcbpractice}

\booksubsection{冗餘文件處理之策}

手機用久則“長胖”，多半乃冗餘文件之作祟。

\begin{itemize}
  \item \textbf{七十二時之自動清理}：微信、QQ等社交App所接收之臨時文件（表情包、聊天圖、小視頻），大多觀一次即無用。於 [微信] > [設置] > [通用] > [存儲空間] 中開"自動清理"，將過期緩存設爲七十二時。手機頓然清爽，且不丟重要之記錄。
  \item \textbf{存儲空間之定向清理}：安卓 [設置] > [存儲] > [清理加速]，iOS [設置] > [通用] > [iPhone 存儲空間]，系統將按"緩存 / 臨時文件 / 可卸載App / 大文件"分類列出。重點清理：(1) 緩存文件；(2) 已過期之安裝包（升級App後所留之舊版安裝包）；(3) 重複之文件（同名同尺寸之照）。
  \item \textbf{自動分類之規}：諸品牌相冊App支持智能分類：人物、風景、文檔、截圖。開啓後，新拍之合影自動歸入"人物 > 張三"，掃描之合同自動歸入"PDF"。\textbf{配以每月之整理，可大降"尋物不得"之概率}。
\end{itemize}

\booksubsection{圖之安全處理之規}

圖乃手機中最“敏感”之數據類型——一張隨手之拍即能暴露汝之位置、設備、時間。吾等建議\textbf{建“雙重保護”之工作流}：

\begin{tcbwarn}{發送證件照/敏感截圖之標準流程}
\begin{enumerate}
  \item \textbf{可見信息打碼}：用系統自帶或第三方之具（如美圖秀秀、Snapseed）\textbf{塗抹身份證號之中間數位、車牌號、門牌、姓名、金額}等關鍵字段；
  \item \textbf{元數據之剝離}：發送前於相冊中點"分享"或"導出"，系統將問是否保留位置之信息，擇\textbf{"不保留"}。如此拍攝之設備型號、GPS座標、拍攝之時間即不隨圖同發；
  \item \textbf{臨時水印}（可選）：于敏感截圖上加"僅供XX辦理使用"之文字水印，避被濫用。
\end{enumerate}

\end{tcbwarn}

元數據（EXIF）清理之常用方法（基礎操作，深度防護見第六章“照片脫敏處理”）：

\begin{itemize}
  \item \textbf{系統自帶}：iOS [設置] > [隱私與安全性] > [相機]，可閉"精確位置"；分享時擇"僅當前圖"將\textbf{自動剝離大部EXIF}。
  \item \textbf{安卓}：諸品牌相冊App于"分享"時將提示"是否移除位置之信息"，勾選可也。
  \item \textbf{專業之具}：ExifTool（電腦端）、Metapho（iOS）、Photo Metadata Remover（安卓）可徹底清除所有元數據。
\end{itemize}

對於需\textbf{長期保存}之敏感圖（身份證、銀行卡、戶口本之掃描件、護照），建議存於\textbf{加密相冊}：

\begin{itemize}
  \item iPhone：[照片] > 選中照片 > 隱藏（僅Face ID驗證後可覽）；
  \item iPhone：[照片] > 選中照片 > 隱藏（僅Face ID驗證後可覽）；
  \item 第三方加密App：\textbf{360安全相冊}、\textbf{騰訊相冊管家（隱私相冊）}、\textbf{KeepSafe}等，單獨密碼保護，隱蔽性更佳。
\end{itemize}
